\section{Introduction}\label{sec:introduction}

A path that crosses a wide strip must be long, but its critical weight is small. The quantitative question is how these two effects balance. For the honeycomb lattice we prove that a strip of height $h$ has bridge mass comparable to $h^{-1/4}$, while the first moment of bridge length has mass at most $Ch^{13/12}$. A fixed fraction of the bridge mass has length at least $ch^{4/3}$. The strip mass follows from a finite transfer calculation. To control the length moments we also prove a bound for polygons through two specified edges on a cylinder, uniform even when the marked separation is much smaller than its period.

We realize the honeycomb lattice by joining the centers of adjacent triangles in the unit equilateral triangular tiling. A path in a tiled domain begins and ends at midpoints of boundary edges, enters and leaves normally, and visits each triangle center at most once. Its length $\ell(\gamma)$ is the number of centers visited. Thus a path with $\ell$ centers has $\ell-1$ full edges between centers and two terminal half-edges. Throughout the paper its weight is
\[
 w(\gamma)=\kappa^{\ell(\gamma)},\qquad
 \kappa=(2\cos(\pi/8))^{-1}.
\]
For polygons, length is the number of centers, equivalently the number of edges. Polygons are simple, unoriented, and unrooted unless specified otherwise. Sums with specified marks count each polygon once.

Let $d_0=\sqrt3/2$ be the spacing between successive rows of triangles. A \emph{pure lattice cut} is a line bounding triangular rows, parallel to one of the three edge directions of the tiling; a pure-sided tiled domain has sides on such cuts. Fix one boundary midpoint of the infinite strip between two parallel lattice cuts at distance $h d_0$, where $h\ge1$ is an integer. Let $\mathcal B_h$ be the paths from this point to any midpoint on the opposite cut, with all visited centers strictly between the cuts. Put
\[
 B_h=\sum_{\gamma\in\mathcal B_h}w(\gamma),\qquad
 B_h(E)=\sum_{\gamma\in\mathcal B_h}\mathbf1_E(\gamma)w(\gamma).
\]
The lattice symmetries make these quantities independent of the selected midpoint and of the pure lattice orientation of the strip. Diameters refer to Euclidean distance in this fixed embedding, including the terminal midpoints; changing to the center trace changes diameter by a bounded amount.

\begin{theorem}[Finite bridge estimates]\label{thm:finite-bridges}
There are constants $c,C>0$ such that, for every sufficiently large integer $h$,
\begin{align}
 c h^{-1/4}&\le B_h\le C h^{-1/4},\label{eq:bridge-mass}\\
 \sum_{\gamma\in\mathcal B_h}\ell(\gamma)w(\gamma)&\le C h^{13/12},\label{eq:bridge-first-moment}\\
 B_h\{\ell\ge c h^{4/3}\}&\ge c h^{-1/4}.\label{eq:bridge-long-fraction}
\end{align}
Moreover, for every $h\ge1$ and $x\ge0$,
\begin{equation}\label{eq:bridge-diameter}
 B_h\{\diam\gamma>x\}\le C(1+x)^{-1/4}.
\end{equation}
\end{theorem}

Dividing \eqref{eq:bridge-first-moment} by \eqref{eq:bridge-mass} gives a conditional mean at most $Ch^{4/3}$. The positive mass in \eqref{eq:bridge-long-fraction} is uniform in the height; it does not average over heights or prescribe the terminal port. These finite statements are also useful as inputs to renewal arguments, but their proofs here use finite paths and polygons only.

The marked estimate has additional geometric restrictions. Let $e_1,e_2$ be unit lattice vectors at angle $60^\circ$, and let
\[
 P=P_1e_1+P_2e_2,\qquad P_1,P_2\in\mathbb Z_{>0},\qquad d=P_1+P_2.
\]
Choose a tiling edge in either axial direction $e_1$ or $e_2$, and its translate by $P$. Passing through such a mark means that the polygon uses the honeycomb link crossing that tiling edge. For an integer $k\ge2$, quotient the plane by translation by $kP$ and set $N=kd$. The two marks remain distinct in this cylinder.

\begin{theorem}[Two axial marks]\label{thm:two-marks}\label{u:two-edge-theorem}
If $d$ is even, the critical mass of simple polygons in the cylinder $\mathbb R^2/(kP\mathbb Z)$ through both marks satisfies
\begin{equation}\label{eq:marked-uniform}
 \sum_{\gamma\text{ through both marks}}\kappa^{\ell(\gamma)}
 \le C N^{-4/3}(N/d)^{15/8},
\end{equation}
with one constant $C$ independent of $P_1,P_2,k$ and the axial edge choice.
\end{theorem}

The ratio power is needed when summing over pairs of marks. It is strictly smaller than the dimension two, so the sum over short marked separations converges after dyadic decomposition. We also prove a positive nesting estimate. Let $n_0,\ldots,n_5$ be the six unit normals to the pure lattice directions, in clockwise order with $n_0$ pointing vertically upward. For a vertex $y$ of the triangular tiling and $R\in d_0\mathbb Z_{>0}$, put
\[
 H(y;R)=\{z:n_j\cdot(z-y)\le R\text{ for }0\le j\le5\}.
\]
This is a regular tiled hexagon with actual Euclidean inradius $R$. These are all regular hexagons whose sides are pure cuts; the lattice restriction on $R$ and on the center is verified in Section~\ref{sec:hexagon-completion}.

For a finite tiled domain $D$ and a triangular lattice vertex $y$, let $Z_D(y)$ count systems of mutually disjoint polygons contained in $D$, every polygon surrounding $y$, with weight $\kappa^\ell$ per polygon and a factor $2$ for each polygon. The empty system contributes $1$. Define $Z_2(y;R)=Z_{H(y;R)}(y)$.

\begin{theorem}[Nesting and free boundary length]\label{thm:nesting}
Uniformly over triangular lattice vertices $y$ and $R\in d_0\mathbb Z_{>0}$ with $R\ge1$,
\[
 Z_2(y;R)\asymp R^{1/12}.
\]
More generally, for any fixed $0<a\le b<\infty$, every finite tiled domain satisfying
\[
 B(y,aR)\subset D\subset B(y,bR),\qquad R\ge1,
\]
has $Z_D(y)\asymp_{a,b}R^{1/12}$. In particular this comparison is uniform over all translations and pure-cut roundings satisfying these inclusions.

In $H(y;R)$, give paths between ordered free boundary ports critical weight and condition on diameter at least $R/2$. Their total mass is comparable to $R^{3/4}$, and their mean number of visited centers is comparable to $R^{4/3}$.
\end{theorem}

\subsection{Context and method}

Nienhuis's analysis of the dilute $O(n)$ model predicted the honeycomb critical point and critical exponents \cite{Nienhuis1982}. Duminil-Copin and Smirnov proved the exact connective constant $\sqrt{2+\sqrt2}$ by a parafermionic observable \cite{DuminilCopinSmirnov2012}. Their convention also uses mid-edge endpoints and counts visited vertices. We use the corresponding critical activity and give the local flux calculation with our boundary-port convention.

For critical strip-crossing mass, Beaton, Bousquet-M\'elou, de Gier, Duminil-Copin and Guttmann proved convergence to zero \cite[Theorem~10]{BeatonEtAl2014}. Glazman and Manolescu obtained a logarithmic bound along a subsequence \cite[Proposition~1.1]{GlazmanManolescu2017}, and Krachun and Panagiotis proved a polynomial upper bound with exponent $10^{-10}$ \cite{KrachunPanagiotis2026}. The predicted decay is $h^{-1/4}$, as recorded in \cite{GlazmanManolescu2017}. Theorem~\ref{thm:finite-bridges} gives this decay with fixed multiplicative constants, together with first-moment and positive-mass length estimates under the same finite bridge measure.

Yang--Baxter deformations provide a second context: Glazman and Manolescu proved invariance of the boundary two-point function for columnwise rhombic half-planes with Yang--Baxter weights and angles in $[\pi/3,2\pi/3]$ \cite[Theorem~1]{GlazmanManolescu2017}. The dilute transfer-matrix setting also underlies the bulk spectral calculations of Zhou and Batchelor \cite{ZhouBatchelor1997}. Polynomial exchange and fusion methods for dense and open dilute $O(1)$ models \cite{DiFrancescoZinnJustin2005,GarbaliNienhuis2017Ground,GarbaliNienhuis2017Sum} provide precedents for our finite construction; the periodic $n=0$ identities and their physical contractions are proved here.

For systems of polygons separating two fixed points on a cylinder, we introduce a generating function with a variable factor per polygon. Polynomial exchange and fusion express it as a finite Laurent polynomial divided by a vacuum normalization. We estimate the numerator by a current expansion based on the classical level-one boson construction of Frenkel--Kac and Segal \cite{FrenkelKac1980,Segal1981,Frenkel1985}; the thermal trace estimates and contour deformations are established here.

The denominator may be small, and an unequal marked interval leads to an open product of operators rather than a periodic trace. Our argument treats these two issues separately. At zero separating-loop weight the physical partition is exactly one. This special value supplies a saturation point for the analytic estimate and eventually a bounded-factor normalization. For the unequal marked intervals we construct an operator for the ordered particle expansion, prove compactness with quantitative state bounds, and control products on the two long intervals. A bound for closed periodic traces is used only after the operator realizing the open products has been constructed.

The geometry has an analogous normalization issue. Joining each chord to an exterior connector would appear to lose a critical bridge factor at every stage. We vary the intermediate cuts and apply a second moment estimate. The difference between two cut choices produces additional bridges whose averaged costs recover those apparent losses. A further gain near corners makes the sum over endpoint bins finite. Together with \eqref{eq:marked-uniform}, this gives a second-moment bound at the required scale for the chord length and then the long-bridge mass in \eqref{eq:bridge-long-fraction}.

\subsection{Proof route}

Sections~\ref{c6:strip-transfer} and~\ref{c6:strip-asymptotics} establish the strip mass independently, first through a finite double-row interpolation formula and then through positive interpolation estimates. Section~\ref{c6:periodic-polynomial} constructs the periodic vacuum vectors, which represent the empty-connectivity eigenvalue, and identifies their physical polygon pairing. Section~\ref{c6:gas} transforms that finite polynomial into a current expansion, an algebraic reorganization of its particle sum, and bounds it, including all zero modes and contour changes. Section~\ref{c6:subheading-09} obtains the preliminary saturation needed to identify a nonzero physical normalization.

Section~\ref{c6:geometry} proves the boundary crossing and cylinder-barrier estimates, and converts positive nesting into a free-boundary length sum. Section~\ref{sec:uniform-nesting} then sharpens the normalization to bounded factors. Section~\ref{c6:subheading-14} inserts prescribed-edge observables, called magnetic insertions in the transfer calculation, and proves the unequal marked estimate. Finally, Section~\ref{sec:sewing} closes and extends chords in the exterior, sums over marked pairs, and proves Theorem~\ref{thm:finite-bridges}. Each stage keeps the physical length convention above; no change to a length-conditioned or infinite-path law occurs in the argument.
